% Pravila Boolove algebre in pravila za kvantifikatorje iz poglavja Boolova
% algebra na listu A4.
\documentclass[10pt]{article}

\usepackage[a4paper, landscape, margin=10mm]{geometry}
\usepackage[slovene]{babel}

\usepackage{amsmath}
\usepackage{mathtools}

\usepackage{fontspec}
\usepackage{unicode-math}
\setmainfont{Libertinus Serif}[Ligatures = TeX, Numbers = OldStyle]
\setmathfont{Libertinus Math}
\newcommand{\lmnVerzalke}{\addfontfeature{Numbers=Lining}}
\usepackage[protrusion = true, expansion = true, final]{microtype}

\usepackage{xcolor}
\definecolor{ULantracit}{HTML}{58595B}
\definecolor{ULsiva}{HTML}{A7A8AA}
\colorlet{barvaPasu}{ULsiva!35}

\usepackage{multicol}
\setlength{\columnsep}{6mm}

\usepackage{eqparbox}

\newcommand{\nls}[1]{\guillemotright\textit{#1}\guillemotleft} % naravni jezik
\newcommand{\df}[1]{\emph{\textbf{#1}}}  % definiranec
\newcommand{\defeq}{\mathrel{{:}{=}}}  % definicijska enacba
\newcommand{\dfeq}{\defeq}                     % starejse ime

\newcommand{\konst}[1]{\mathrm{k}_{#1}} % konstantna preslikava


% Substitucija
\newcommand{\xsubst}[2]{#1[#2]}
\newcommand{\subst}[2]{(#1)[#2]}
\newcommand{\subto}{\,{\mapsto}\,}

%%% quantifiers
\newcommand{\all}[1]{\forall #1 \,.\,}
\newcommand{\some}[1]{\exists #1 \,.\,}
\newcommand{\exactlyone}[1]{\exists! #1 \,.\,}
\newcommand{\descr}[1]{\iota #1 \,.\,}

\newcommand{\defiff}{\quad {{:}{\liff}}\quad }

% \Box je v unicode-math prazen kvadratek; rabimo ga kot vsebnik.
\renewcommand{\Box}{\mdlgwhtsquare}

\newcommand{\brac}[1]{\ensuremath{\cdots}\enspace$\langle$#1$\rangle$\enspace\ensuremath{\cdots}}

%\newcommand{\lam}[3]{\lambda #1 \,{\in}\, #2\,.\left(#3\right)}

% \newcommand{\uall}[2]{\forall\, #1\,.\left(#2\right)}
% \newcommand{\usome}[2]{\exists\, #1\,.\left(#2\right)}
% \newcommand{\uexactlyone}[3]{\exists!\, #1\,.\left(#2\right)}
% \newcommand{\ulam}[2]{\lambda #1 .\left(#2\right)}

% \newcommand{\xall}[3]{\forall\, #1 \,{\in}\, #2\,.\,#3}
% \newcommand{\xsome}[3]{\exists\, #1 \,{\in}\, #2\,.\,#3}
% \newcommand{\xexactlyone}[3]{\exists!\, #1 \,{\in}\, #2\,.\,#3}
% \newcommand{\xlam}[3]{\lambda #1 \,{\in}\, #2\,.\,#3}

% \newcommand{\xuall}[2]{\forall\, #1\,.\,#2}
% \newcommand{\xusome}[2]{\exists\, #1\,.\,#2}
% \newcommand{\xuexactlyone}[2]{\exists!\, #1,.\,#2}
% \newcommand{\xulam}[2]{\lambda #1 .\,#2}

% \newcommand{\tlam}[3]{\lambda #1 \,{:}\, #2\,.\,\left(#3\right)}
% \newcommand{\xtlam}[3]{\lambda #1 \,{:}\, #2\,.\,#3}

\newcommand{\iinn}{\qquad\text{in}\qquad}

% Implikacija
\newcommand{\lthen}{\Rightarrow}
\newcommand{\liff}{\Leftrightarrow}

% Hipoteticna izpeljava: iz predpostavk #1 pridemo do #2.
\NewDocumentCommand{\izpeljava}{m m}{\begin{array}[b]{c}{#1}\\ \vdots\\ {#2}\end{array}}
\newcommand{\hypoinfer}{\izpeljava}  % starejse ime

% Standardne množice
\newcommand{\NN}{\mathbb{N}}
\newcommand{\ZZ}{\mathbb{Z}}
\newcommand{\QQ}{\mathbb{Q}}
\newcommand{\RR}{\mathbb{R}}
\newcommand{\CC}{\mathbb{C}}

\newcommand{\one}{\mathbb{1}}
\newcommand{\two}{\mathbb{2}}
\newcommand{\unit}{(\,)}

%%% Množice
\newcommand{\set}[1]{\left\{#1\right\}}
\newcommand{\such}{\; \middle| \;}
\newcommand{\iso}{\cong}

%%%%%%  Projections and Injections
%%%%%%%%%%%%%%%%%%%%%%%%%%%%%%%%%%%%%%%%%%%%%%%%%%%%%%%%%%%%%
\newcommand{\inj}[1][]{\mathsf{in}_{#1}}
\newcommand{\pr}[1][]{\mathsf{pr}_{#1}}

%% Neobvezna nadpisa loci vejica, kadar sta podana oba.
\NewDocumentCommand{\nadpisa}{m m}{#1\IfBlankF{#2}{,}#2}

\NewDocumentCommand{\fst}{O{} O{}}{\pr[1]^{\nadpisa{#1}{#2}}}
\NewDocumentCommand{\snd}{O{} O{}}{\pr[2]^{\nadpisa{#1}{#2}}}
\NewDocumentCommand{\inl}{O{} O{}}{\inj[1]^{\nadpisa{#1}{#2}}}
\NewDocumentCommand{\inr}{O{} O{}}{\inj[2]^{\nadpisa{#1}{#2}}}


%%% Naravna števila
\newcommand{\suc}[1]{{#1^{+}}}

\newcommand{\Set}{\mathsf{Set}}
\newcommand{\On}{\mathsf{On}}

%% Drevesa
\newcommand{\Tree}{\mathsf{Tree}}
\newcommand{\tree}[1]{\mathsf{tree}(#1)}
\newcommand{\emptyTree}{\mathsf{empty}}

%%% Operacije na množicah

\newcommand{\pow}[1]{\mathcal{P}(#1)}       % potenčna množica
\newcommand{\npow}[2]{\mathcal{P}_{#1}(#2)} % podmnožice dane velikosti

%% Preslikave
\newcommand{\id}[1][]{\mathrm{id}_{#1}}
\newcommand{\inv}[1]{{#1}^{-1}}
\newcommand{\restrict}[2]{#1\mathrel{\big|}_{#2}}

\newcommand{\diag}[1][]{\Delta_{#1}}

\newcommand{\img}[1]{#1_{*}}
\newcommand{\invimg}[1]{#1^{*}}

%% Kombinatorika
\newcommand{\ceil}[1]{\lceil #1 \rceil}

%% Relacije
\newcommand{\transpose}[1]{#1^{\mathsf{T}}}
\newcommand{\compl}[1]{#1^{\mathsf{C}}}
\newcommand{\rel}[1]{\mathrel{#1}}


\pagestyle{empty}
\setlength{\parindent}{0pt}

%% Razdelek: pas z imenom in znakom, nato pravila.
\newlength{\razmikRazdelkov} \setlength{\razmikRazdelkov}{2ex}

\newcommand{\pas}[2]{%
  {\setlength{\fboxsep}{1.5pt}%
   \colorbox{barvaPasu}{\makebox[\dimexpr\linewidth - 2\fboxsep\relax][l]{%
       \strut\hspace{1pt}\textbf{#1}\hfill\ensuremath{#2}\hspace{1pt}}}}}

\newenvironment{razdelek}[2]{%
  \par\noindent\begin{minipage}{\linewidth}%
    \pas{#1}{#2}\par\vspace{1ex}}
 {\end{minipage}\par\vspace{\razmikRazdelkov}}

\newcommand{\opomba}[1]{\par\vspace{0.3ex}{\footnotesize\itshape #1}\par}

%% Pravilo: leve strani iz iste skupine so enako široke, zato so relacije
%% na isti navpičnici. Ime pravila stoji ob desnem robu.
\newlength{\razmikPravil} \setlength{\razmikPravil}{1ex}
\newcommand{\skupina}{izjave}
\newcommand{\relacija}{=}
\newcommand{\pravilo}[3][]{%
  \par\eqmakebox[\skupina][r]{$#2$}${}\relacija #3$%
  \hfill{\footnotesize\color{ULantracit}\textsc{#1}}\par\vspace{\razmikPravil}}

\begin{document}

\raggedcolumns
\begin{multicols}{3}

{\Large\bfseries Pravila Boolove algebre}\par
\vspace{0.3ex}
{\footnotesize\color{ULantracit}\textsc{Logika in množice}}\par
\vspace{\razmikRazdelkov}

\begin{razdelek}{Konjunkcija}{\land}
  \pravilo[asociativnost $\land$]{(p \land q) \land r}{p \land (q \land r)}
  \pravilo[komutativnost $\land$]{p \land q}{q \land p}
  \pravilo[idempotentnost $\land$]{p \land p}{p}
  \pravilo[$\top$ je nevtralen za $\land$]{\top \land p}{p}
  \pravilo[$\bot$ absorbira $\land$]{\bot \land p}{\bot}
\end{razdelek}

\begin{razdelek}{Disjunkcija}{\lor}
  \pravilo[asociativnost $\lor$]{(p \lor q) \lor r}{p \lor (q \lor r)}
  \pravilo[komutativnost $\lor$]{p \lor q}{q \lor p}
  \pravilo[idempotentnost $\lor$]{p \lor p}{p}
  \pravilo[$\bot$ je nevtralen za $\lor$]{\bot \lor p}{p}
  \pravilo[$\top$ absorbira $\lor$]{\top \lor p}{\top}
\end{razdelek}

\begin{razdelek}{Implikacija}{\lthen}
  \pravilo[kontrapozicija]{(p \lthen q)}{(\neg q \lthen \neg p)}
  \pravilo{(p \lthen q)}{\neg p \lor q}
  \pravilo{(\bot \lthen q)}{\top}
  \pravilo{(\top \lthen q)}{q}
  \pravilo{(p \lthen \bot)}{\neg p}
  \pravilo{(p \lthen \top)}{\top}
  \pravilo{(p \lthen p)}{\top}
  \pravilo{((p \land q) \lthen r)}{(p \lthen (q \lthen r))}
  \pravilo{((p \lor q) \lthen r)}{(p \lthen r) \land (q \lthen r)}
  \pravilo{(p \lthen (q \land r))}{(p \lthen q) \land (p \lthen r)}
\end{razdelek}

\begin{razdelek}{Ekvivalenca}{\liff}
  \pravilo{(p \liff q)}{(p \lthen q) \land (q \lthen p)}
  \pravilo{(p \liff \top)}{p}
  \pravilo{(p \liff \bot)}{\neg p}
\end{razdelek}

\begin{razdelek}{Kombinirana pravila}{}
  \pravilo[De Morganovo pravilo za $\land$]{\neg(p \land q)}{\neg p \lor \neg q}
  \pravilo[De Morganovo pravilo za $\lor$]{\neg(p \lor q)}{\neg p \land \neg q}
  \pravilo{\neg(p \lthen q)}{p \land \neg q}
  \pravilo[absorpcijsko pravilo za $\land$]{p \land (p \lor q)}{p}
  \pravilo[absorpcijsko pravilo za $\lor$]{p \lor (p \land q)}{p}
  \pravilo[distributivnost $\land$]{p \land (q \lor r)}{(p \land q) \lor (p \land r)}
  \pravilo[distributivnost $\lor$]{p \lor (q \land r)}{(p \lor q) \land (p \lor r)}
\end{razdelek}

\begin{razdelek}{Negacija}{\lnot}
  \pravilo{\neg \top}{\bot}
  \pravilo{\neg \bot}{\top}
  \pravilo[negacija je involucija]{\neg\neg p}{p}
  \pravilo[pravilo o izključenem tretjem]{p \lor \neg p}{\top}
  \pravilo{p \land \neg p}{\bot}
\end{razdelek}

\end{multicols}

\renewcommand{\skupina}{kvantifikatorji}
\renewcommand{\relacija}{\iff}

\begin{multicols}{2}

\begin{razdelek}{Univerzalni kvantifikator}{\forall}
  \pravilo{(\all{x \in \emptyset} \phi(x))}{\top}
  \pravilo{(\all{x \in \set{a}} \phi(x))}{\phi(a)}
  \pravilo{(\neg \all{x \in A} \phi(x))}{\some{x \in A} \neg \phi(x)}
  \pravilo{(\psi \lthen \all{x \in A} \phi(x))}{\all{x \in A} (\psi \lthen \phi(x))}
  \pravilo{(\psi \lor \all{x \in A} \phi(x))}{\all{x \in A} (\psi \lor \phi(x))}
  \pravilo{(\all{x \in A} \phi(x) \land \psi(x))}{(\all{y \in A} \phi(y)) \land (\all{z \in A} \psi(z))}
  \pravilo{(\all{x \in A} \all{y \in B} \phi(x, y))}{\all{y \in B} \all{x \in A} \phi(x, y)}
  \pravilo{(\all{u \in A \times B} \phi(u))}{\all{x \in A} \all{y \in B} \phi(x, y)}
  \pravilo{(\all{u \in A + B} \phi(u))}{(\all{x \in A} \phi(\inl(x))) \land (\all{y \in B} \phi(\inr(y)))}
  \pravilo{(\all{u \in A \cup B} \phi(u))}{(\all{x \in A} \phi(x)) \land (\all{y \in B} \phi(y))}
  \pravilo{(\all{u \in \set{x \in A \such \psi(x)}} \phi(u))}{\all{x \in A} (\psi(x) \lthen \phi(x))}
\end{razdelek}

\begin{razdelek}{Eksistenčni kvantifikator}{\exists}
  \pravilo{(\some{x \in \emptyset} \phi(x))}{\bot}
  \pravilo{(\some{x \in \set{a}} \phi(x))}{\phi(a)}
  \pravilo{(\neg \some{x \in A} \phi(x))}{\all{x \in A} \neg \phi(x)}
  \pravilo{(\some{x \in A} \phi(x)) \lthen \psi}{\all{x \in A} (\phi(x) \lthen \psi)}
  \pravilo{(\psi \land \some{x \in A} \phi(x))}{\some{x \in A} (\psi \land \phi(x))}
  \pravilo{(\some{x \in A} \phi(x) \lor \psi(x))}{(\some{y \in A} \phi(y)) \lor (\some{z \in A} \psi(z))}
  \pravilo{(\some{x \in A} \some{y \in B} \phi(x, y))}{\some{y \in B} \some{x \in A} \phi(x, y)}
  \pravilo{(\some{u \in A \times B} \phi(u))}{\some{x \in A} \some{y \in B} \phi(x, y)}
  \pravilo{(\some{u \in A + B} \phi(u))}{(\some{x \in A} \phi(\inl(x))) \lor (\some{y \in B} \phi(\inr(y)))}
  \pravilo{(\some{u \in A \cup B} \phi(u))}{(\some{x \in A} \phi(x)) \lor (\some{y \in B} \phi(y))}
  \pravilo{(\some{u \in \set{x \in A \such \psi(x)}} \phi(u))}{\some{x \in A} (\psi(x) \land \phi(x))}
\end{razdelek}

\end{multicols}

\opomba{V pravilih, v katerih nastopa $\psi$ brez argumenta, spremenljivka $x$ v~$\psi$ ne nastopa.}

\end{document}

%%% Local Variables:
%%% mode: latex
%%% TeX-engine: luatex
%%% End:
